\documentclass[aspectratio=169]{beamer}
\usetheme{Madrid}
\usecolortheme{beaver}
\usepackage[utf8]{inputenc}
\usepackage[T1]{fontenc}
\usepackage{tikz}
\usepackage{amssymb}
\usepackage{graphicx}
\setbeamertemplate{navigation symbols}{}
\setbeamercolor{title}{fg=white,bg=blue!55!black}
\setbeamercolor{frametitle}{fg=white,bg=blue!45!black}
\setbeamercolor{block title}{bg=blue!15,fg=black}
\setbeamercolor{block body}{bg=blue!5,fg=black}

\title{M4 — Lists}
\subtitle{Creating, indexing, appending, looping — list + loop = superpower}
\author{ByteForge --- Grade 8 --- App Builder}
\date{}

\begin{document}

\begin{frame}
    \titlepage
\end{frame}

\begin{frame}{Today's Plan}
    \begin{itemize}
        \item A List of Values
        \item Tuples
        \item sorted() and Slicing
        \item Filtering
        \item The "in" Keyword
        \item[$\bigstar$] BOSS --- M4 Capstone
    \end{itemize}
\end{frame}

\begin{frame}{A list is an ordered collection of values}
    \Large
    A list holds many values under one name, in order — \texttt{for temp in temperatures:} visits each one in turn, exactly like Grade 7's array ideas.
\end{frame}


\begin{frame}{Tuples — a fixed little group of values}
    \Large
    \texttt{(r, g, b)} groups a FIXED number of related values as one item. \texttt{r, g, b = color} splits one tuple into three named variables in a single line.
\end{frame}


\begin{frame}{sorted() and slicing — built-in tools, not loops}
    \Large
    \texttt{sorted(player\_scores)} returns a NEW list in order — no loop required. \texttt{[-5:]} means "the last 5 items", a negative index counting from the end.
\end{frame}


\begin{frame}{Filtering — building a new, smaller list}
    \Large
    Looping through a list and building a NEW list containing only items that pass a test — "give me only the results that match X" is an everywhere-task in real software.
\end{frame}

\begin{frame}{Real-World Example}
    \Large
    Every "sort by price" or "filter by category" button on a shopping site is exactly \texttt{sorted()} and filtering, at real scale. Minecraft's own inventory sorting (a real feature in some versions) is literally \texttt{sorted()} applied to your item list.
\end{frame}


\begin{frame}{Let Them Discover}
    \Large
    Find a real app with a "filter" feature (a shopping site, a music app). What single condition do you think its filter checks per item — the same "loop + condition + build a new list" pattern from this chapter?
\end{frame}

\begin{frame}{Your Turn}
    \Large
    Open the student map, find \textbf{M4 — Lists}, and work through it
    hands-on --- everything you just saw, they now do themselves.
\end{frame}

\end{document}
